% GENERATED by tools/build.sh from main.aux: labels of main.tex, prefixed M-
\newlabel{M-sec:related}{{2}{2}{Related Work}{section.2}{}}
\newlabel{M-fig:pipeline}{{1}{2}{Conditional memory utility evaluation. \textbf {(a)}~Page-visit memory units extracted from past episodes are stored as Raw, Summary or Structured representations without exposing the writer to future evaluation tasks. \textbf {(b)}~For a memory--task pair, a frozen reader is evaluated with identical filler $F$, either with or without the target memory $m$; the paired reward difference defines $\Delta U(m,t)$ under the recorded evaluation conditions. \textbf {(c)}~Repeating this comparison across tasks, representations and replicates yields a task-conditioned utility profile. Formation, storage, retrieval and exposure costs are recorded alongside the measurements (the cost ledger)}{figure.caption.10}{}}
\newlabel{M-sec:method}{{3}{2}{Conditional Memory Evaluation}{section.3}{}}
\newlabel{M-eq:du}{{1}{3}{The measured quantity}{equation.3.1}{}}
\newlabel{M-sec:screen}{{3.2}{3}{The answer-bearing screen is an instrument, not a retriever}{subsection.3.2}{}}
\newlabel{M-sec:metrics}{{3.4}{4}{Noise, labels and cost}{subsection.3.4}{}}
\newlabel{M-sec:setup}{{4}{4}{Benchmarks and Setup}{section.4}{}}
\newlabel{M-tab:coverage}{{1}{4}{Every run a reported number comes from. Counts are memory--task pairs times the conditions each pair was evaluated under (representations, arms or replicates, as named). A dataset row is one replicate of one pair under one condition and records both arms. Excluded rows had no producible reward and are dropped, never scored as zero}{table.caption.19}{}}
\newlabel{M-fig:sparse}{{2}{5}{Answer bearing is sparse at every screen cap (exploratory). \textbf {(a)}~Number of memory units (of $\bearUnits $) that bear on each count of screened future tasks, at the paper's cap of eight units per task (log axis). \textbf {(b)}~Share of bearing memories that bear on exactly one task when the cap is 4, 8, 16 or 32 units; $n$ is the number of bearing memories. The median breadth is one at every cap}{figure.caption.20}{}}
\newlabel{M-sec:results}{{5}{5}{Results}{section.5}{}}
\newlabel{M-sec:scarcity}{{5.1}{5}{Specific answers are sparse over memories}{subsection.5.1}{}}
\newlabel{M-fig:robust}{{3}{5}{The result holds (LongMemEval-V2). \textbf {(a)}~Bearing minus control, paired within memory, in the original run, in a new session with the memory's position redrawn per replicate, and under a second reader (95\% intervals over memories). \textbf {(b)}~Bearing-cell $\Delta U$ of each representation at its own length ($n=\lcBearCells $) and with every arm padded to the Raw arm's token count ($n=\soFN $, excluding cells exposed to a cache-key defect); task-clustered intervals. Pre-registered: the original run and the second reader in (a); exploratory: the new session in (a) and panel (b)}{figure.caption.21}{}}
\newlabel{M-sec:screened}{{5.2}{5}{Screening makes the effect measurable}{subsection.5.2}{}}
\newlabel{M-tab:screened}{{2}{6}{Screened run on LongMemEval-V2: the same 60 memories paired with three strata of tasks. Both-arms-fail is the fraction of cells in which the reader fails with \emph {and} without the memory; such cells cannot show a memory effect. The criterion for a material reduction (\textsc {bearing}'s Wilson upper bound below the comparator's point estimate) was fixed in code before these numbers existed}{table.caption.22}{}}
\newlabel{M-sec:similarity}{{5.3}{6}{Retrievers find answer-bearing memories, but miss most}{subsection.5.3}{}}
\newlabel{M-fig:mech}{{4}{6}{What carries utility (LongMemEval-V2). \textbf {(a)}~Bearing-cell $\Delta U$ as Raw is cut to its last $N$ tokens ($\lcBearCells $ cells, task-clustered intervals); labels give the share of cells whose text still contains the answer. \textbf {(b)}~Share of a task's bearing units that a retriever ranks in its top 1, 5 and 10 of the domain pool ($\soScrDisc $ tasks; task-clustered interval shown at top 10); the dashed line is chance. Both panels are exploratory; the matched-length contrast in (a) was pre-registered}{figure.caption.23}{}}
\newlabel{M-sec:repr}{{5.4}{6}{Evaluated compressed representations gain less than Raw}{subsection.5.4}{}}
\newlabel{M-sec:readertransfer}{{5.5}{7}{The contrast keeps its direction under a second reader}{subsection.5.5}{}}
\newlabel{M-sec:r0}{{5.6}{7}{What the instrument did not find}{subsection.5.6}{}}
\newlabel{M-sec:discussion}{{6}{7}{Discussion and Limitations}{section.6}{}}
\newlabel{M-sec:limitations}{{6}{7}{Discussion and Limitations}{section.6}{}}
\newlabel{M-fig:cells}{{5}{8}{What single cells can and cannot show (exploratory). \textbf {(a)}~Per-cell labels (helped, unchanged, hurt) in the original run against a new session at the same memory position ($\soStabSameN $ replicate cells; diagonal agrees). At a new position $\kappa $ falls to $\soStabDiffK $, and between two readers it is $\raReaderKappa $. \textbf {(b)}~Help among cells the reader got wrong without the memory (filled) and harm among cells it got right (hollow), by stratum, Wilson 95\% intervals: help differs sharply, harm does not detectably}{figure.caption.24}{}}
\newlabel{M-fig:position}{{6}{8}{Position affects per-pair measurements; the aggregate bearing--control contrast remains positive after position standardization (exploratory). Per-cell $\Delta U$ by the memory's position among the five context units in the second-session run, with task-clustered 95\% bands: the bearing effect dips in the middle and recovers at the end, while the control stays near zero at every position}{figure.caption.26}{}}
\newlabel{M-sec:release}{{6}{8}{Release and reproducibility}{section*.29}{}}
\newlabel{M-tocindent-1}{0pt}
\newlabel{M-tocindent0}{0pt}
\newlabel{M-tocindent1}{4.185pt}
\newlabel{M-tocindent2}{10.34999pt}
\newlabel{M-tocindent3}{0pt}
\newlabel{M-tocindent4}{0pt}
\newlabel{M-tocindent5}{0pt}
\newlabel{M-TotPages}{{9}{9}{}{page.9}{}}
